% !TeX root = project-report-en-revised.tex
% !TeX program = lualatex
% Build: latexmk -lualatex -interaction=nonstopmode -halt-on-error -outdir=.build/project-report-en-revised project-report-en-revised.tex
\documentclass[11pt,a4paper]{article}
\usepackage[margin=21mm,headheight=14pt]{geometry}
\usepackage{lmodern,amsmath,amssymb,graphicx,booktabs,array,tabularx}
\usepackage{xcolor,enumitem,fancyhdr,caption,xurl,placeins,float,needspace,titlesec}
\usepackage[colorlinks=true,linkcolor=teal,citecolor=teal,urlcolor=teal]{hyperref}
\definecolor{reportblue}{HTML}{18536B}
\hypersetup{pdftitle={Effects of Feature Groups on Four-Lepton MC Event Selection},pdfauthor={HiggsML project}}
\pagestyle{fancy}
\fancyhf{}
\fancyhead[L]{\small\textcolor{reportblue}{HiggsML | Project report}}
\fancyhead[R]{\small Revised 15 November 2026}
\fancyfoot[C]{\thepage}
\setlength{\parindent}{0pt}
\setlength{\parskip}{6pt}
\setlength{\emergencystretch}{2em}
\setlist[itemize]{leftmargin=1.4em,itemsep=3pt,topsep=4pt}
\captionsetup{font=small,labelfont=bf,justification=raggedright,singlelinecheck=false}
\newcolumntype{Y}{>{\raggedright\arraybackslash}X}
\titleformat{\section}{\Large\bfseries\color{reportblue}}{\thesection}{1em}{}
\titleformat{\subsection}{\large\bfseries\color{reportblue}}{\thesubsection}{1em}{}
\titlespacing*{\section}{0pt}{16pt}{8pt}
\titlespacing*{\subsection}{0pt}{12pt}{5pt}
\raggedbottom
\widowpenalty=10000
\setlength{\textfloatsep}{12pt}
\setlength{\intextsep}{12pt}
\newcommand{\gev}{\,\mathrm{GeV}}

\begin{document}
{\LARGE\bfseries\color{reportblue} Effects of Feature Groups on\\ Four-Lepton MC Event Selection\par}
\vspace{6pt}
Project report $\cdot$ Feature-group deletion comparison $\cdot$ Revised 15 November 2026

\section{Research purpose and comparison scope}

This study examines how kinematic inputs affect the selection of simulated $H\to ZZ^*\to4\ell$ events with a neural network.
One input group is removed at a time, and the network is retrained using the same samples, hidden-layer architecture and threshold rule.
The comparison reports signal yield retention, background yield retention and the signal-to-background ratio after selection. Mass spectra are also examined for selected models.
The resulting differences describe changes in the selection procedure after retraining and threshold selection. They are not interpreted as independent physical contributions from each input group.

The inputs are divided into four groups. A contains the four leptons' transverse momenta and pseudorapidities. B contains the two pair masses and within-pair angular separations. C contains the four-lepton transverse momentum and the azimuthal difference between the pairs. D contains five production and decay angles.
All groups are obtained from the same lepton four-momenta, so their information can overlap.
For example, pair masses and single-lepton momenta may provide correlated information for separating signal from background.

The comparison includes the full input set ABCD and four reduced configurations: BCD, ACD, ABD and ABC omit A, B, C and D, respectively.
Each configuration is separately initialized and trained with seeds 42, 43 and 44, giving fifteen models.
All models use the same event split. The validation sample is used to select each checkpoint and score threshold. Final-evaluation results are calculated after both are frozen.
Differences are calculated within each seed, then summarized across the three seeds.

In the current three-seed comparison, removing C decreases signal retention, increases background retention and decreases S/B in every seed.
For the other deletions, the direction depends on the metric and seed; see Section~\ref{sec:results}.
This agreement in direction is a descriptive finding, not evidence of statistical significance.
The ABCD and BCD mass spectra also show that the relative background contribution in each mass region can change as total background falls.

The report uses controlled-MC results from the draft-v2 three-subset analysis. Appendix~\ref{app:provenance} identifies the source and the scope of earlier checks.
Numerical values are retained from the existing evidence record. The evaluation and report artifacts bound to that record are missing from the current working directory. Appendix~\ref{app:availability} explains how this limits the present numerical checks.
This is an educational and technical study, not a measurement using real collision data.

\section{Simulated samples and event selection}

\subsection{Signal, background and study population}
The samples are taken from the ATLAS 2020 educational 13 TeV four-lepton collection~\cite{OpenData}.
The release description and file list are available directly in \href{https://opendata.cern.ch/record/15005}{CERN Open Data record 15005}.
Only two MC files are used: signal DSID 345060 (\texttt{ggH125\_ZZ4lep}) and background DSID 363490 (\texttt{llll}).
The signal models gluon-fusion production of a Higgs boson with mass set to $125\gev$, followed by $H\to ZZ^*\to4\ell$. The background sample also contains four-lepton events.
Process names follow the sample identifiers. Detailed production settings remain unverified; see Appendix~\ref{app:fields}.
Throughout this report, background refers only to the selected \texttt{llll} sample. The quoted S/B therefore does not represent the full background composition of an experimental analysis.

Both samples undergo the same physical selection. It requires an event trigger, exactly four leptons passing the quality cuts, same-flavour opposite-sign pairs, and the specified momentum, acceptance and pair-mass cuts.
The analysis is restricted to two electrons and two muons ($2e2\mu$) with $105\leq m_{4\ell}<140$ GeV.
The electron and muon pairs form two candidates. The pair whose mass is closer to $91.1876$ GeV is denoted $Z_1$ and the other $Z_2$.
Appendix~\ref{app:cuts} lists the full requirements in execution order and explains their relation to the cited analyses. Appendix~\ref{app:fields} defines the four-momenta and angles.

In the final-evaluation sample, ``before network selection'' refers to all events that pass the common physical cuts.
``After network selection'' means the high-score subset that also passes the frozen network-score threshold.

\subsection{Sample allocation and subset roles}
After physical selection, event groups are assigned to training, validation and final evaluation with probabilities of 60\%, 20\% and 20\%.
The realized subset sizes need not match these percentages exactly.
Each physical event and its related records remain in one subset. All input configurations and training seeds use the same assignment.
The training sample is used to fit standardization and network parameters. Validation selects the checkpoint and threshold, and final evaluation gives the yields and spectra for the frozen selection.

\begin{center}
\small
\begin{tabularx}{\linewidth}{@{}p{.18\linewidth}p{.19\linewidth}Y@{}}
\toprule
Subset & Allocation probability & Role \\
\midrule
Training & 60\% & Fit input-standardization and network parameters. \\
Validation & 20\% & Select a checkpoint using AUC, then determine a threshold from background scores. \\
Final evaluation & 20\% & Calculate selected yields and mass spectra after freezing the model and threshold. \\
\bottomrule
\end{tabularx}
\end{center}
The checkpoint is selected using validation AUC with absolute weights. The selected network then supplies the validation-background scores used to set the threshold.
Final-evaluation events do not enter parameter updates, checkpoint selection or threshold determination.

\begin{table}[H]
\caption{Downloaded file entries and final-evaluation events after physical selection.}
\label{tab:population}
\centering\small
\begin{tabularx}{\linewidth}{@{}Y p{.23\linewidth} rrr@{}}
\toprule
Sample stage & Counting unit & Signal & Background & Total \\
\midrule
Downloaded MC files & File entries & 164,716 & 554,279 & 718,995 \\
Final evaluation (before selection) & Distinct physical events & 9,049 & 698 & 9,747 \\
\bottomrule
\end{tabularx}
\end{table}
The rows of Table~\ref{tab:population} describe different processing stages.
The second includes both physical selection and subset allocation and cannot be obtained by multiplying the first by 20\%.
The available aggregates do not provide complete group counts for training and validation. These counts are not inferred from the allocation probabilities.

\subsection{Event counts, physical weights and expected yields}\label{sec:weights}
Event counts give the size of the simulated sample. Expected yields give the number of events predicted after physical normalization.
Each physical event is counted once; yields instead sum the signed physical weights of its simulation records.
Both samples are normalized to $10\,\mathrm{fb}^{-1}$, the project's MC normalization target.

For generator weight $w_i^{\rm gen}=\texttt{mcWeight}$, the physical weight of record $i$ is
\[
 w_i^{\rm phys}=\frac{\mathcal L\,\sigma_{\rm eff}}{S_{\rm gen}}\,w_i^{\rm gen},
 \qquad \mathcal L=10^4\,\mathrm{pb}^{-1}.
\]
Here $\mathcal L$ is the target integrated luminosity and $\sigma_{\rm eff}$ the effective sample cross section. The official parameter table supplies the generator-weight sum $S_{\rm gen}$.
The pinned official table~\cite{SampleInfo} agrees with the project configuration:
\begin{center}
\small
\begin{tabular}{lrr}
\toprule
Sample & $\sigma_{\rm eff}$ ($\mathrm{pb}$) & $S_{\rm gen}$ \\
\midrule
Signal & 0.0060239 & 27881776.6536 \\
Background & 1.2578 & 7538705.8077 \\
\bottomrule
\end{tabular}
\end{center}

Final evaluation has allocation probability $20\%=0.2$.
To estimate the yield after common physical selection at the same target luminosity, the evaluation weight is
\[
 u_i=\frac{w_i^{\rm phys}}{0.2}.
\]
Let $E$ contain the final-evaluation events of one process that pass the physical cuts. Denote the network score by $s_i$ and the validation-background threshold by $t$, as defined in Section~\ref{sec:threshold}.
The before-selection and high-score yields are
\[
 Y_{\rm before}=\sum_{i\in E}u_i,\qquad
 Y_{\rm high}=\sum_{i\in E,\,s_i\geq t}u_i.
\]
$Y_{\rm before}$ and $Y_{\rm high}$ are the corresponding expected yields.
Weight signs are retained, so simulation-record counts cannot be substituted for expected yields.
Other normalization settings and differences from the official example appear in Appendix~\ref{app:normalization}.
\FloatBarrier

\section{Feature groups and neural-network method}

\subsection{Input information and learning objective}
The inputs are kinematic quantities from leptons passing the common physical selection.
Table~\ref{tab:feature-overview} summarizes the four groups. The transverse momentum $p_T$ is perpendicular to the beam, pseudorapidity $\eta$ describes direction relative to the beam, and $\Delta R$ measures separation in the pseudorapidity--azimuth plane.
All inputs are obtained from the released lepton four-momenta.
Appendix~\ref{app:fields} gives variable names, reconstruction formulas and angular conventions.

\begin{table}[H]
\caption{Feature groups and dimensions. ABCD, BCD, ACD, ABD and ABC have 19, 11, 15, 17 and 14 inputs respectively.}
\label{tab:feature-overview}
\centering\small
\begin{tabularx}{\linewidth}{@{}cp{.36\linewidth}Yc@{}}
\toprule
Group & Variables & Physical meaning & Dim. \\
\midrule
A & Four leptons' $p_T$, $\eta$ & Single-lepton momentum scale and direction & 8 \\
B & $m_{Z1}$, $m_{Z2}$, within-pair $\Delta R$ & Pair masses and internal angular separation & 4 \\
C & $p_{T,4\ell}$, $|\Delta\phi_{ZZ}|$ & Overall transverse motion and azimuthal relation between pairs & 2 \\
D & Three angle cosines and two plane angles & Production and decay orientations & 5 \\
\bottomrule
\end{tabularx}
\end{table}
Each configuration is trained separately.
No classifier takes $m_{4\ell}$ as an explicit input. The other variables are constructed from lepton four-momenta, however, and can remain correlated with mass.

Signal and background have training labels 1 and 0.
Each input is standardized using statistics from the training sample. For event $i$ and input $j$, let $x_{ij}$ be the original value. Its training mean and standard deviation are $\mu_j$ and $\sigma_j$:
\[
 \widetilde x_{ij}=\frac{x_{ij}-\mu_j}{\sigma_j}.
\]
Standardization puts the inputs on comparable numerical scales.
Validation and evaluation use the same training-derived $\mu_j$ and $\sigma_j$.
If the training standard deviation is zero, the denominator is set to 1.

\subsection{Fully connected model with three hidden layers}\label{sec:network}
The classifier is a PyTorch multilayer perceptron (MLP) with three hidden layers:
\[
 \underbrace{d_{\rm in}}_{\text{input}}
 \;\longrightarrow\;
 \underbrace{64\;\longrightarrow\;64\;\longrightarrow\;32}_{\text{three hidden layers}}
 \;\longrightarrow\;
 \underbrace{1}_{\text{output}}.
\]
The input dimension $d_{\rm in}$ depends on the configuration in Table~\ref{tab:feature-overview}.
All configurations share the hidden-layer architecture; Table~\ref{tab:network} lists the operations.

\begin{table}[H]
\caption{Network computation order. Linear denotes a fully connected linear transformation. Configurations change the input dimension while retaining the same hidden-layer architecture.}
\label{tab:network}
\centering\small
\begin{tabularx}{\linewidth}{@{}p{.20\linewidth}p{.18\linewidth}Y@{}}
\toprule
Component & Dimensions & Operations in order \\
\midrule
Input processing & $d_{\rm in}\to d_{\rm in}$ & Standardize each input using its training-set mean and standard deviation. \\
Hidden layer 1 & $d_{\rm in}\to64$ & Linear, LayerNorm, SiLU, Dropout (0.1). \\
Hidden layer 2 & $64\to64$ & Linear, LayerNorm, SiLU, Dropout (0.1). \\
Hidden layer 3 & $64\to32$ & Linear, LayerNorm, SiLU. \\
Output & $32\to1$ & Linear produces a real value $z$; sigmoid is applied when computing the score. \\
\bottomrule
\end{tabularx}
\end{table}
A node in a fully connected layer computes
\[
 a_j=\sum_k W_{jk}h_k+b_j.
\]
Here $h_k$ is an output from the previous layer; $W_{jk}$ and $b_j$ are the connection weights and biases.
Hidden layers use LayerNorm and SiLU; the first two also use Dropout with probability 0.1.
Appendix~\ref{app:layers} defines these operations and explains how they differ between training and evaluation.

\Needspace{7\baselineskip}
The output layer produces a real value $z$, converted to a classification score by sigmoid:
\[
 s=\frac{1}{1+e^{-z}}
\]
A higher score ranks an event as more signal-like.
The optimization weights are rescaled within each class. The resulting score cannot be read directly as a true signal fraction or a calibrated signal probability.

\subsection{Training and configuration comparison}
Training minimizes mini-batch weighted binary cross entropy.
Each epoch shuffles the training events and divides them into batches of at most 1024 events. The network computes a score $s_i$ for each event.
With labels $y_i=1$ for signal and $y_i=0$ for background, the per-event loss is
\[
 \ell_i=-\left[y_i\ln s_i+(1-y_i)\ln(1-s_i)\right].
\]
For signal and background this becomes $\ell_i=-\ln s_i$ and $\ell_i=-\ln(1-s_i)$ respectively.
Minimizing these terms favours higher signal scores and lower background scores.

Optimization weights use absolute physical weights divided by the mean within each training class:
\[
 \overline w_c=\frac{1}{N_c}\sum_{i\in\mathcal T_c}|w_i^{\rm phys}|,
 \qquad
 \alpha_i=\frac{|w_i^{\rm phys}|}{\overline w_{y_i}}.
\]
Here $c=0,1$ denotes background or signal, $\mathcal T_c$ contains that class's training events and $N_c$ is its event count.
The loss weights are therefore nonnegative, and the mean per-event optimization weight in each class is 1. Relative weights within a class are preserved.
Class counts may differ, so total optimization weights need not be equal.

For a batch $\mathcal B$ of $m$ events, the objective is
\[
 L_{\rm batch}=\frac{1}{m}\sum_{i\in\mathcal B}\alpha_i\ell_i.
\]
Backpropagation calculates its parameter gradients, and AdamW updates parameters with learning rate $10^{-3}$ and weight decay $10^{-4}$.
One pass through all training events is an epoch. Each epoch ends with validation, followed by the checkpoint and early-stopping rules below.
Optimization weights are used only for learning. Final yields use the signed physical weights in Section~\ref{sec:weights}.

All five configurations share the optimizer, hidden-layer architecture, epoch limit, event allocation and checkpoint rule.
Each is initialized and trained with seeds 42, 43 and 44.
The seed controls initialization, shuffling and Dropout. Changing the input dimension also changes the number of parameters in the first layer.
Each reduced classifier is thus fitted separately under the same training rules. An input group is not simply set to zero in the full network during evaluation.

\subsection{Checkpoint selection with validation AUC}
After each epoch, AUC is calculated on validation events using absolute physical weights.
The best checkpoint is updated only when AUC improves by more than $10^{-4}$ over the saved best value.
Training stops after 20 epochs without such improvement or at 200 epochs, and the saved best parameters are restored.
ABCD seed 42 stopped at epoch 87 and uses epoch 67; BCD stopped at epoch 86 and uses epoch 66.
The score threshold is determined after checkpoint selection.
Appendix~\ref{app:network} gives implementation and reproduction details.

\subsection{Threshold selection with validation background}\label{sec:threshold}
For each model, the threshold $t$ is the median of its validation-background scores, weighted by absolute physical weights.
Scores are ordered from low to high and absolute weights are accumulated.
The threshold is the observed score at the first position reaching half the total weight, without interpolation.
Events with $s\geq t$ enter the high-score category; those with $s<t$ enter the low-score category. Events with tied scores remain together.
The threshold is fixed before final evaluation and is not adjusted using evaluation results.

ABCD and BCD seed-42 thresholds are $0.746724$ and $0.860891$ respectively.
Score scales can differ across networks, so threshold magnitude is not itself a performance measure.
The comparison applies each model's own threshold to the common evaluation events and compares the resulting yields.

\subsection{Application to final evaluation after freezing}
Input-standardization parameters, the network checkpoint and the threshold are frozen before final evaluation.
Evaluation events are standardized and scored using the saved parameters, then assigned to the high- or low-score category. Each category is summarized by physical-group counts, signed yields and mass spectra.
Network selection changes which events are retained, not their individual masses.

Signal retention, background retention and S/B use the high-score yields over the full $[105,140)$ GeV window.
The sum of high- and low-score yields is checked against the total before selection.
The validation-background weighted median does not require exactly 50\% of the final-evaluation background yield to be retained. Appendix~\ref{app:threshold} explains why and gives an illustrative example.

\section{Comparison metrics and worked example}\label{sec:basic}

\subsection{Event counts and weighted yields}
Table~\ref{tab:groups} uses ABCD and BCD seed 42 to show physical-group counts in the common evaluation sample and each high-score subset.
ABCD uses all inputs; BCD removes the single-lepton group A.
For example, ABCD retains 8,761 of 9,049 signal groups.
These counts describe the size of the finite simulation sample. Expected yields instead sum signed weights, as defined in Section~\ref{sec:weights}.

\begin{table}[H]
\caption{Physical event groups in final evaluation. The last two columns are the seed-42 high-score subsets.}
\label{tab:groups}
\centering
\begin{tabular}{lrrr}
\toprule
Process & Before selection & ABCD high score & BCD high score \\
\midrule
Signal & 9,049 & 8,761 & 8,563 \\
Background & 698 & 386 & 378 \\
\bottomrule
\end{tabular}
\end{table}

\subsection{Retention fractions and signal-to-background ratio}
Selection is described by signal retention, background retention and S/B.
Let $E_S$ and $E_B$ contain the signal and background events in final evaluation before network selection. Their weights $u_i$ are defined in Section~\ref{sec:weights}.
Before-selection yields are
\[
 S_0=\sum_{i\in E_S}u_i,\qquad
 B_0=\sum_{i\in E_B}u_i.
\]
For score $s_i$ and the model's frozen threshold $t$, selected yields are
\[
 S=\sum_{i\in E_S,\,s_i\geq t}u_i,\qquad
 B=\sum_{i\in E_B,\,s_i\geq t}u_i.
\]
All four yields correspond to $10\,\mathrm{fb}^{-1}$; the metrics are
\[
 R_S=\frac{S}{S_0},\qquad R_B=\frac{B}{B_0},\qquad Q=\frac{S}{B}.
\]
$R_S$ and $R_B$ are signed-yield retention fractions, and background rejection is $1-R_B$.
$Q$ is the selected S/B; for example, $Q=2$ means the signal yield is twice the selected background yield.
Retention is displayed as a percentage; yields and S/B are reported directly.

Substituting $S=R_SS_0$ and $B=R_BB_0$ gives
\[
 Q=\frac{S_0}{B_0}\frac{R_S}{R_B}.
\]
With common $S_0$ and $B_0$, the two retention fractions determine S/B. The three metrics therefore do not provide independent evidence.
$Q$ can rise when background falls proportionally more, even while signal is lost.
The metrics are reported together to describe this tradeoff. Configurations are not ranked by $Q$ alone, and the metrics are not combined into a total score.
S/B does not account for the full counting and systematic uncertainty model. An increase in S/B alone cannot establish improved statistical sensitivity or measurement precision.

All models use the same threshold rule, but their evaluation background retention need not be equal.
The comparison uses the yields obtained with each fixed threshold. Thresholds are not readjusted to equalize signal or background retention.

If $S_0$ or $B_0$ is nonpositive, the corresponding retention fraction is not calculated.
If $B\leq0$, S/B is not calculated and the reason is stated.
Yields always retain the signs of the weights. Negative yields and retention fractions outside $[0,1]$ are reported and flagged because they limit interpretation.

\subsection{ABCD and BCD worked example}
\begin{table}[H]
\caption{Expected signed yields and retention at $10\,\mathrm{fb}^{-1}$. Errors are grouped MC standard errors at fixed networks and thresholds. Retention errors include covariance from shared before/after events.}
\label{tab:yields}
\centering
\begin{tabular}{llrr}
\toprule
Process & Selection & Expected yield & Retention (\%) \\
\midrule
Signal & Before selection ($S_0$) & $2.7710\pm0.0293$ & --- \\
 & ABCD high score & $2.6824\pm0.0288$ & $96.80\pm0.19$ \\
 & BCD high score & $2.6215\pm0.0285$ & $94.60\pm0.24$ \\
Background & Before selection ($B_0$) & $2.3553\pm0.1683$ & --- \\
 & ABCD high score & $1.2027\pm0.1202$ & $51.06\pm3.57$ \\
 & BCD high score & $1.1012\pm0.1156$ & $46.75\pm3.57$ \\
\bottomrule
\end{tabular}
\end{table}
Applying the metric definitions to the yields in Table~\ref{tab:yields} gives Table~\ref{tab:worked-metrics}.
The last column is BCD minus ABCD, showing changes after removing A.

\begin{table}[H]
\caption{Seed-42 calculations and changes after removing A. The displayed calculations use rounded values; complete results were calculated from unrounded aggregates. Retention differences are in percentage points (pp), while S/B differences are direct differences.}
\label{tab:worked-metrics}
\centering\small
\renewcommand{\arraystretch}{2.6}
\begin{tabular}{lccr}
\toprule
Metric & ABCD & BCD & BCD $-$ ABCD \\
\midrule
Signal retention $R_S$ & $\dfrac{2.6824}{2.7710}\simeq96.80\%$ & $\dfrac{2.6215}{2.7710}\simeq94.60\%$ & $-2.20$ pp \\
Background retention $R_B$ & $\dfrac{1.2027}{2.3553}\simeq51.06\%$ & $\dfrac{1.1012}{2.3553}\simeq46.75\%$ & $-4.31$ pp \\
S/B $Q$ & $\dfrac{2.6824}{1.2027}\simeq2.23$ & $\dfrac{2.6215}{1.1012}\simeq2.38$ & approx. $+0.15$ \\
\bottomrule
\end{tabular}
\end{table}
BCD seed 42 retains less signal and less background. Background retention falls proportionally more, so S/B rises.
This is a selection tradeoff; one seed cannot establish that BCD is generally superior.

\Needspace{9\baselineskip}
\textbf{Finite-MC uncertainty.}
The values following $\pm$ in Table~\ref{tab:yields} are standard errors from finite MC samples.
For a given process and selection, let $\mathcal G$ contain the contributing physical groups, indexed by $g$. The records contributing to each group are $\mathcal I_g$, with signed evaluation weights $u_i$.
Different physical groups are treated as independent counting units under a Poisson-type MC variance convention:
\[
 U_g=\sum_{i\in\mathcal I_g}u_i,\qquad
 Y=\sum_{g\in\mathcal G}U_g,\qquad
 \sigma_Y=\sqrt{\sum_{g\in\mathcal G}U_g^2}.
\]
Records from the same physical event are combined before calculating errors, so correlated records are not treated as independent events.
The error reduces to $\sqrt{N}$ only in special cases such as independent unit-weight counting.
These errors are conditional on fixed networks, thresholds and normalization parameters. They exclude retraining, threshold estimation and systematic uncertainties in cross sections, luminosity or scale factors.
Retention numerators and denominators share events, so their covariance is included as described in Appendix~\ref{app:errors}.

\subsection{Within-seed deletion contrasts and aggregation}
Removing A, B, C and D gives BCD, ACD, ABD and ABC respectively.
For each training seed, the contrast subtracts ABCD from the reduced configuration.
For removed group $G$ and seed $r$,
\[
 \Delta R_{S,G,r}=R_S(\mathrm{ABCD}\setminus G,r)-R_S(\mathrm{ABCD},r),
\]
\[
 \Delta R_{B,G,r}=R_B(\mathrm{ABCD}\setminus G,r)-R_B(\mathrm{ABCD},r),\qquad
 \Delta Q_{G,r}=Q(\mathrm{ABCD}\setminus G,r)-Q(\mathrm{ABCD},r).
\]
All contrasts follow reduced minus ABCD.
Retention differences are multiplied by 100 and reported in percentage points; $\Delta Q$ is a direct S/B difference.

A negative $\Delta R_S$ means less signal retained, a negative $\Delta R_B$ less background retained and a positive $\Delta Q$ higher S/B.
The signs must therefore be interpreted jointly with signal loss and background suppression.

For each deletion and metric, the median and minimum--maximum range include the differences from all three seeds.
Each subtraction compares configurations with the same training-seed label. Subtracting configuration medians calculated separately is generally not equivalent.

If all three differences have the same sign, these three training results agree in direction. Mixed signs indicate that the direction depends on the seed.
The range is not a confidence interval and cannot establish statistical significance.
Both the reduced network and its threshold are refitted. The differences therefore describe the whole selection method, not an independent physical contribution from the removed group.
\FloatBarrier

\section{Results for the five input configurations}\label{sec:results}

\subsection{Retention and S/B for individual seeds}
Table~\ref{tab:basic-all} lists all fifteen models.
The existing evidence record gives common before-selection yields of $S_0=2.7710$ and $B_0=2.3553$. For each model, high- and low-score yields sum to the total before selection.
Differences in the table therefore arise from the events retained by each model.
All seeds are reported. The following subsection summarizes the configuration contrasts and their variation across seeds.

\begin{table}[H]
\caption{Yields, retention and S/B for five configurations and three seeds. Values retain the existing record's fifteen valid model summaries and calculations from unrounded yields; Appendix~\ref{app:availability} limits the current checks. $R_S$ and $R_B$ are percentages; $S$ and $B$ correspond to $10\,\mathrm{fb}^{-1}$.}
\label{tab:basic-all}
\centering\small
\begin{tabular}{llrrrrr}
\toprule
Inputs & Seed & $S$ & $B$ & $R_S$ & $R_B$ & $Q=S/B$ \\
\midrule
ABCD & 42 & 2.6824 & 1.2027 & 96.80 & 51.06 & 2.230 \\
ABCD & 43 & 2.6809 & 1.1589 & 96.75 & 49.20 & 2.313 \\
ABCD & 44 & 2.6670 & 1.0834 & 96.25 & 46.00 & 2.462 \\
\addlinespace
BCD & 42 & 2.6215 & 1.1012 & 94.60 & 46.75 & 2.381 \\
BCD & 43 & 2.6156 & 1.1145 & 94.39 & 47.32 & 2.347 \\
BCD & 44 & 2.6264 & 1.1344 & 94.78 & 48.16 & 2.315 \\
\addlinespace
ACD & 42 & 2.6667 & 1.1536 & 96.24 & 48.98 & 2.312 \\
ACD & 43 & 2.6679 & 1.1232 & 96.28 & 47.69 & 2.375 \\
ACD & 44 & 2.7039 & 1.2813 & 97.58 & 54.40 & 2.110 \\
\addlinespace
ABD & 42 & 2.5731 & 1.2186 & 92.86 & 51.74 & 2.112 \\
ABD & 43 & 2.5408 & 1.2192 & 91.69 & 51.76 & 2.084 \\
ABD & 44 & 2.5805 & 1.2026 & 93.13 & 51.06 & 2.146 \\
\addlinespace
ABC & 42 & 2.6975 & 1.1241 & 97.35 & 47.73 & 2.400 \\
ABC & 43 & 2.6972 & 1.0311 & 97.33 & 43.78 & 2.616 \\
ABC & 44 & 2.6941 & 1.1386 & 97.22 & 48.34 & 2.366 \\
\bottomrule
\end{tabular}
\end{table}

\subsection{Deletion contrasts and seed variation}
Table~\ref{tab:basic-deltas} applies Section~\ref{sec:basic} to form reduced-minus-ABCD contrasts.
Figure~\ref{fig:basic-deltas} shows the same differences relative to zero.
The table ranges and connecting lines are descriptive summaries; three same-sign values are not evidence of statistical significance.

\begin{table}[H]
\caption{Deletion contrasts relative to ABCD within the same seed. Retention differences are in percentage points; S/B differences are dimensionless. Every summary requires all three seeds.}
\label{tab:basic-deltas}
\centering\small
\begin{tabular}{llccccc}
\toprule
Removed & Metric & Seed 42 & Seed 43 & Seed 44 & Median & Min--max \\
\midrule
A & $\Delta R_S$ & $-2.20$ & $-2.35$ & $-1.47$ & $-2.20$ & $[-2.35,-1.47]$ \\
A & $\Delta R_B$ & $-4.31$ & $-1.88$ & $+2.17$ & $-1.88$ & $[-4.31,+2.17]$ \\
A & $\Delta Q$ & $+0.150$ & $+0.033$ & $-0.147$ & $+0.033$ & $[-0.147,+0.150]$ \\
\addlinespace
B & $\Delta R_S$ & $-0.57$ & $-0.47$ & $+1.33$ & $-0.47$ & $[-0.57,+1.33]$ \\
B & $\Delta R_B$ & $-2.09$ & $-1.51$ & $+8.40$ & $-1.51$ & $[-2.09,+8.40]$ \\
B & $\Delta Q$ & $+0.081$ & $+0.062$ & $-0.351$ & $+0.062$ & $[-0.351,+0.081]$ \\
\addlinespace
C & $\Delta R_S$ & $-3.94$ & $-5.05$ & $-3.12$ & $-3.94$ & $[-5.05,-3.12]$ \\
C & $\Delta R_B$ & $+0.68$ & $+2.56$ & $+5.06$ & $+2.56$ & $[+0.68,+5.06]$ \\
C & $\Delta Q$ & $-0.119$ & $-0.229$ & $-0.316$ & $-0.229$ & $[-0.316,-0.119]$ \\
\addlinespace
D & $\Delta R_S$ & $+0.54$ & $+0.59$ & $+0.98$ & $+0.59$ & $[+0.54,+0.98]$ \\
D & $\Delta R_B$ & $-3.34$ & $-5.42$ & $+2.34$ & $-3.34$ & $[-5.42,+2.34]$ \\
D & $\Delta Q$ & $+0.169$ & $+0.302$ & $-0.096$ & $+0.169$ & $[-0.096,+0.302]$ \\
\bottomrule
\end{tabular}
\end{table}

\begin{figure}[H]
\centering
\includegraphics[width=\linewidth]{figures/basic-20261108/deletion-differences.pdf}
\caption{Within-seed deletion differences relative to ABCD. The horizontal axis gives the removed group. Panels show signal-retention differences (pp), background-retention differences (pp) and S/B differences. Circles denote seeds 42--44, open diamonds the medians and thin lines the minimum--maximum ranges. Dashed lines mark zero. The ranges are not confidence intervals.}
\label{fig:basic-deltas}
\end{figure}

\subsection{Interpretation of the results}
\textbf{Removing C.}
Signal retention decreases in all three seeds ($-5.05$ to $-3.12$ pp), background retention increases ($+0.68$ to $+5.06$ pp) and S/B decreases ($-0.316$ to $-0.119$).
For each frozen selection, ABCD retains more signal and less background than ABD. Both retention metrics therefore have more favourable point estimates for ABCD.
C contains the four-lepton transverse momentum and the azimuthal difference between pairs.
The two variables were not removed separately, so their individual roles cannot be determined. The comparison also cannot separate correlations and refitting effects into independent contributions.

\textbf{Removing A.}
Signal retention decreases consistently ($-2.35$ to $-1.47$ pp), whereas the directions of background retention and S/B depend on the seed.
Seeds 42 and 43 reduce both signal and background and increase S/B; seed 44 reduces signal, increases background and lowers S/B.
Signal loss is consistent across these seeds, but the extent to which background suppression offsets it depends on the trained model.

\textbf{Removing B.}
The deletion differences change sign across seeds for all three metrics.
Seeds 42 and 43 slightly reduce signal and background retention and increase S/B; seed 44 retains more signal and background and lowers S/B.
A median alone would conceal these changes in direction.
The results do not support characterizing B as consistently beneficial or consistently redundant.

\textbf{Removing D.}
Signal retention increases in every seed ($+0.54$ to $+0.98$ pp), but background retention and S/B do not improve consistently.
Seeds 42 and 43 reduce background and increase S/B, while seed 44 reverses those directions.
Higher signal retention after removing D does not establish a consistent improvement in overall selection. It also does not rule out the usefulness of angles with other input combinations or analysis conditions.

Of the four deletions, only removing C reduces signal retention and increases background retention in all three seeds.
For every other deletion, at least one metric changes direction across seeds. The results therefore do not give a consistent ordering of the configurations.
These contrasts are conditional on the current MC sample, network and threshold rules. The report does not estimate their uncertainties, including correlations from shared evaluation events.
Directional agreement across seeds cannot establish significance.
The results describe selection tradeoffs. They do not identify universally optimal inputs or establish a gain in measurement precision.
\FloatBarrier

\section{Changes in mass spectra}
The ABCD and BCD seed-42 mass spectra show how selection changes the mass distributions as well as the total yields.
Absolute spectra show the expected yield in each region. Separately normalized shapes distinguish a reduction in total yield from a change in relative regional contributions.
The example is limited to these two models. It does not explain all deletion results or identify the mechanism behind the C-group result.
Plots use seven 5 GeV mass bins, while retention and S/B use the full mass window.

\subsection{Expected yields in mass bins}
Signal and background are separately binned in four-lepton invariant mass.
For mass bin $b$,
\[
 Y_b=\sum_{i\in b}u_i.
\]
Before selection the sum includes all evaluation events; afterwards it includes only high-score events.
Bin sums equal the corresponding total yield.
Figure~\ref{fig:yields} places signal on the left and background on the right, with common luminosity and bin edges before and after selection.

\begin{figure}[H]
\centering
\includegraphics[width=.98\linewidth]{figures/basic-20261108/mass-yields.pdf}
\caption{Signed mass spectra in final evaluation. ABCD is above BCD, both for seed 42; signal is on the left and background on the right. Blue circles show before-selection yields and orange circles high-score yields. Each point is a 5 GeV bin; lines guide the eye and do not define a continuous mass model. The vertical axis gives expected yields at $10\,\mathrm{fb}^{-1}$, with different signal/background scales. Error bars are grouped MC standard errors. Shading flags limited statistical support, and negative yields are retained.}
\label{fig:yields}
\end{figure}
The simulated signal is concentrated near 125 GeV, and both high-score selections retain this main region.
The network only decides whether events remain; it does not change their masses.
These coarse bins cannot establish a shifted peak, improved detector resolution or a changed intrinsic Higgs width.

Before network selection, the background extends nonuniformly across the full mass window.
Regional yields change by different amounts, with a more pronounced reduction in the high-mass tail.
The next subsection describes these changes through each region's share of the total.
\FloatBarrier

\subsection{Normalized mass-spectrum shapes}
Figure~\ref{fig:shapes} normalizes each curve separately to compare regional shares.
For a process and selection, let $Y_{\rm total}=\sum_bY_b$.
When the total is positive,
\[
 f_b=\frac{Y_b}{Y_{\rm total}}.
\]
$f_b$ is the bin's share of that curve's total.
The curve before selection and the high-score curve are each divided by their own total.
The implementation additionally requires both totals to pass the support conditions in Appendix~\ref{app:shape-support}.
Signed bin shares can be negative, so the normalized curve need not be a probability density.
Table~\ref{tab:fractions} sums the shares over specified regions and reports percentages.

\begin{figure}[H]
\centering
\includegraphics[width=.96\linewidth]{figures/basic-20261108/mass-shapes.pdf}
\caption{Separately normalized before/after shapes using the same events and thresholds as Figure~\ref{fig:yields}. Each bin is divided by its curve's own positive signed total. Error propagation includes the shared normalization; before/after curves also share events. Signed normalized shapes need not be probability densities. Shaded and negative bins are retained without smoothing, clipping or rebinning.}
\label{fig:shapes}
\end{figure}

\begin{table}[H]
\caption{Signed yield in fixed mass regions as a percentage of each curve's own total. These are descriptive point estimates without displayed uncertainties and cannot establish significant before/after differences.}
\label{tab:fractions}
\centering\small
\begin{tabular}{llrrr}
\toprule
Process & Region (GeV) & Before selection & ABCD high score & BCD high score \\
\midrule
Signal & $[120,130)$ & 87.9 & 88.6 & 88.5 \\
Background & $[120,130)$ & 32.5 & 51.8 & 47.9 \\
Background & $[135,140)$ & 22.0 & 6.5 & 9.0 \\
\bottomrule
\end{tabular}
\end{table}
The signal share in $[120,130)$ GeV changes little.
The background share rises in that region and falls in $[135,140)$ GeV.
Because total background also decreases, a larger regional share does not imply a larger absolute regional yield.
Uncertainties on the differences between these fractions are not shown, so the changes are interpreted descriptively.

The shape changes are consistent with mass-dependent selection: excluding explicit $m_{4\ell}$ does not guarantee mass independence.
This example cannot attribute the changes to one input or establish a physical background resonance.
\FloatBarrier

\subsection{Statistical support and interpretation limits}
Shading identifies bins with too few groups, unstable yield estimates or strong cancellation between positive and negative weights. Appendix~\ref{app:errors} gives the criteria.
ABCD high-score background has a negative yield in $[105,110)$ GeV.
Several background bins and the signal high-mass tail also have limited support.
Values are shown as recorded, without clipping negative yields to zero.

Error bars describe finite-MC standard errors at fixed networks and thresholds.
The curves before and after selection can contain the same events. Normalized bins also share a common total.
Points and curves are therefore correlated and cannot be treated as independent measurements.

\clearpage
\section{Conclusions and limitations}\label{sec:conclusion}
The report compares the full input set ABCD with four single-group deletions using the same controlled-MC samples and training rules. Each configuration has three training seeds.
Signal retention, background retention and S/B are evaluated at frozen thresholds. Mass spectra for ABCD and BCD describe the accompanying changes in the mass distributions.

Removing C lowers signal retention, raises background retention and lowers S/B in all three seeds.
For deletions of A, B and D, at least one metric changes direction across seeds. These results do not give a consistent feature-group ordering.
The ABCD/BCD examples also show that a region can receive a larger share of the remaining background while total background falls.
Signal retention, background suppression and mass distributions should therefore be considered together. An increase in S/B alone cannot establish better measurement precision.

These conclusions apply only to the two selected samples and the specified network and threshold rules.
The background does not cover the full experimental composition.
Inputs can be correlated, and both the network and threshold are refitted after deletion. The contrasts therefore describe changes in the complete procedure, not independent physical contributions.
The spectra cover only two models at one seed and do not explain the mechanism behind the C-group result.

Minimum--maximum ranges across seeds describe variation in the observed training results. MC standard errors are conditional on fixed networks and thresholds.
Neither includes the full training and threshold-estimation uncertainty, and statistical significance of deletion differences has not been established.
Further evaluation would need to account for shared events when estimating contrast uncertainties, as well as variation from training and threshold estimation. It would also require validation with a specified background composition and physically applicable normalization.
These are not completed results of this report.

The analysis remains exploratory and non-independent.
Repartitioning or rerunning previously used source events does not create an independent validation sample.
Production settings, release preselection, external support for some cuts and physical-weight applicability remain unresolved.
The missing result package also limits end-to-end numerical checks; see Appendix~\ref{app:availability}.
The project is educational and technical, not an official ATLAS/CMS result, a Higgs discovery or a physics measurement.

\clearpage
\appendix

\section{Fields, reconstruction and input conventions}\label{app:fields}

\subsection{MC event generation and released samples}\label{app:mc-generation}
Monte Carlo event generation samples kinematic configurations from theoretical differential cross sections.
Consider idealized independent sampling from a nonnegative distribution, with kinematics $x$ and integrated cross section $\sigma$:
\[
 p(x)=\frac{1}{\sigma}\frac{d\sigma}{dx},\qquad
 \int g(x)p(x)\,dx\simeq\frac{1}{N}\sum_{k=1}^{N}g(x_k),\quad x_k\sim p.
\]
$g$ is the event function being averaged and $N$ the number of simulated events.
This illustrates how random samples approximate distribution integrals; it does not require the actual generator to use equal-weight sampling.
Practical calculations may use importance sampling and event weights, and higher-order calculations and shower matching can produce signed weights.
Physical yields therefore use the weighted sums in Section~\ref{sec:weights}, rather than unit-weight record counts.
The presence of negative weights does not identify their particular generation mechanism.

A typical hadron-collider simulation first generates hard-scattering events using parton distribution functions and perturbative calculations. It then simulates parton showers, hadronization and unstable-particle decays, followed by detector and trigger response and object reconstruction.
The ATLAS simulation-framework description explains how event generators are coupled to detector simulation~\cite{ATLASSimulation}.
That description covers the general simulation chain. It does not document the full production settings of these two 2020 educational samples.

The project uses MC files already generated and released by ATLAS, beginning with their reconstructed lepton fields.
It does not generate hard scattering, showers or detector response.
Here ``reconstruction'' means combining lepton four-momenta, pairing candidates and calculating kinematic variables. These operations use the released reconstructed objects, not generator-truth momenta.
Common physical cuts, event-group allocation, network training and frozen evaluation follow.
These downstream steps are distinct from upstream MC production.

\subsection{Evidence available for sample production}
The inspected CERN release record~\cite{OpenData} identifies the 2020 13 TeV collection of events with at least four leptons (electrons or muons). It also describes loose object and event preselection before release.
The pinned sample table~\cite{SampleInfo} provides DSIDs, names and normalization constants.
Both local ROOT files match the dataset contract in SHA-256, \texttt{mini} tree name and record count. These identity checks do not establish the missing production settings.

\begin{center}
\small
\begin{tabularx}{\linewidth}{@{}>{\raggedright\arraybackslash}p{.23\linewidth}Y@{}}
\toprule
Information & Scope and limitations of the source checks \\
\midrule
Sample identity & Signal DSID 345060 (\texttt{ggH125\_ZZ4lep}) and background DSID 363490 (\texttt{llll}) correspond to the 2020 dataset contract and public file list. \\
Generator and theory settings & DSID-specific production records are still needed for the generator/version, perturbative order, parton distribution functions, shower/hadronization settings, tune and generation filters. \\
Detector and object definitions & The inspected release record, parameter table and ROOT metadata do not establish the detector-simulation configuration, pileup, radiation treatment or exact lepton-object definitions. \\
Release preselection & Loose preselection is confirmed; exact momentum, acceptance and quality thresholds remain unverified. \\
Normalization and background composition & Parameter values can be checked; effective-cross-section factorization, negative-weight origin and the initial-state/interference composition of \texttt{llll} remain unresolved. \\
\bottomrule
\end{tabularx}
\end{center}
DSID 363490 is therefore not described as a verified pure quark-initiated $ZZ$ sample. Generator information from 2025 filenames is not assigned to these 2020 samples.
Downloaded and evaluation counts in Table~\ref{tab:population} cannot determine acceptance relative to all generated events.
Source details, code comparisons and check scope are recorded in \nolinkurl{paper/draft-v2/MC-SOURCE-METHOD-AUDIT-20261115.md}.

\subsection{Read fields}
The \texttt{mini} tree stores events as records: an entry is a record and a branch is a field.
The following table lists all 19 branches read by the project.
The number of branches is distinct from the dimension of the derived network inputs. The source files may contain other branches.

The first seven fields are scalars; the others are per-lepton vectors.
Equal indices across vectors identify the same lepton.
For example, \texttt{lep\_n=4} gives four elements in each lepton vector.

\begingroup
\small
\renewcommand{\arraystretch}{1.13}
\begin{tabularx}{\linewidth}{@{}p{.53\linewidth}Y@{}}
\toprule
Read branches & Meaning and stored units \\
\midrule
\texttt{runNumber, eventNumber, channelNumber} & Event/sample identifiers (integers) \\
\texttt{mcWeight; trigE, trigM; lep\_n} & Generator weight; trigger booleans; lepton count \\
\texttt{lep\_pt, lep\_eta, lep\_phi, lep\_E} & $p_T$, $\eta$, $\phi$, energy; momentum and energy in MeV, angles in radians \\
\texttt{lep\_charge, lep\_type} & Charge and particle type (integer vectors) \\
\texttt{lep\_isTightID, lep\_trigMatched} & Identification and trigger-matching flags (boolean vectors) \\
\texttt{lep\_ptcone30, lep\_etcone20} & Track/calorimeter isolation (MeV) \\
\texttt{lep\_z0; lep\_tracksigd0pvunbiased} & $z_0$ (mm); $d_0/\sigma(d_0)$ (dimensionless) \\
\bottomrule
\end{tabularx}
\endgroup

Kinematic, isolation and impact-parameter fields are floating-point vectors.
Quantities stored in MeV are converted to GeV before use.
Event weights are converted to expected yields as in Section~\ref{sec:weights}.
These field definitions follow the input configuration and earlier ROOT metadata checks. They do not establish the missing production-chain details.

\subsection{Four-momentum reconstruction and input grouping}
\textbf{Single-lepton four-momenta.}
For leptons passing quality selection, momenta and energies are converted to GeV. The leptons are then indexed in descending $p_T$ order.
With $c=1$, released $(p_T,\eta,\phi,E)$ give
\[
 p^\mu=(E,p_x,p_y,p_z),\qquad
 p_x=p_T\cos\phi,\quad p_y=p_T\sin\phi,\quad p_z=p_T\sinh\eta.
\]
These define energy and three momentum components; four-vector sums then provide composite masses and momenta.

\textbf{Pair definitions.}
Same-flavour opposite-sign (SFOS) pairs are formed from electron--positron or oppositely charged muon pairs.
The $2e2\mu$ channel provides a unique electron and muon pair.
The pair closer to $91.1876\gev$ is $Z_1$ and the other $Z_2$.
The reference notation $m_{12},m_{34}$ corresponds to the code fields \texttt{mZ1,mZ2}. The mass criterion defines $Z_1$; it need not be the electron pair and is not ordered by $p_T$.

\textbf{Composite kinematics.}
For pair four-momenta $P_j=p_{j^-}+p_{j^+}$ and total $Q=P_1+P_2$,
\[
 m_{Zj}=\sqrt{P_j^2},\qquad m_{4\ell}=\sqrt{Q^2},\qquad
 p_{T,4\ell}=\sqrt{Q_x^2+Q_y^2}.
\]
The first two quantities are pair and four-lepton masses; the third is the total momentum transverse to the beam.
This combines existing objects without resimulating the detector.
Radiation, detector response and selection still affect the spectra, whose widths are not directly the intrinsic Higgs width.

\textbf{Input groups.}
A takes sorted $p_T,\eta$ directly; B and C use the above four-vectors.
Within-pair $\Delta R$ uses lepton $\eta,\phi$, with azimuthal differences wrapped to $[-\pi,\pi)$.
\texttt{deltaPhi\_ZZ} is the absolute azimuthal difference between pairs.
The 19 inputs below are derived from the read fields, not 19 additional fields already stored in the file.

\begin{table}[H]
\caption{Input groups and physical meanings. Momenta and masses use GeV; azimuthal angles use radians.}
\label{tab:features}
\centering\small
\renewcommand{\arraystretch}{1.12}
\begin{tabularx}{\linewidth}{@{}p{.07\linewidth}p{.43\linewidth}Y@{}}
\toprule
Group & Code variables & Physical meaning \\
\midrule
A & \texttt{lep1\_pt}--\texttt{lep4\_pt}; \texttt{lep1\_eta}--\texttt{lep4\_eta} &
Single-lepton transverse momenta and pseudorapidities, describing momentum scales and directions relative to the beam. \\
B & \texttt{mZ1}, \texttt{mZ2}; \texttt{deltaR\_Z1}, \texttt{deltaR\_Z2} &
Two pair masses and the internal separation of each pair, $\Delta R=\sqrt{(\Delta\eta)^2+(\Delta\phi)^2}$. \\
C & \texttt{pt4l}, \texttt{deltaPhi\_ZZ} &
Four-lepton transverse momentum and the absolute azimuthal difference between the two pairs. \\
D & \texttt{cos\_theta\_star}, \texttt{cos\_theta\_1}, \texttt{cos\_theta\_2};
\texttt{phi\_decay\_planes}, \texttt{phi\_production\_plane} &
One production-angle cosine, two decay polar-angle cosines, and two angles describing the relative orientations of decay and production planes. \\
\bottomrule
\end{tabularx}
\end{table}

\subsection{Five-angle conventions}
Group D compares production and decay directions in specified frames.
Lepton four-momenta are Lorentz-transformed before applying these definitions:
\begin{enumerate}[leftmargin=1.7em,itemsep=3pt]
\item \textbf{Production angle.} In the four-lepton rest frame, compare the $Z_1$ direction with the transformed positive laboratory beam direction to obtain $\cos\theta^*$.
\item \textbf{Two decay polar angles.} In each pair rest frame, compare the negative lepton with the direction opposite the other pair, obtaining $\cos\theta_1,\cos\theta_2$.
\item \textbf{Two plane angles.} In the four-lepton rest frame, a decay-plane normal is the negative-lepton momentum crossed with its positive partner's momentum. The positive beam and $Z_1$ define the production-plane normal. Oriented angles about $Z_1$ compare the two decay planes, and the production plane with the first decay plane.
\end{enumerate}
Angles lie in $[-\pi,\pi)$ and cosines in $[-1,1]$.
Every configuration uses these same conventions.
\FloatBarrier

\section{Complete selection requirements}\label{app:cuts}

\subsection{Order of the common selection}
The table lists physical cuts and their purposes in execution order.
Every event must pass them before network selection, then be allocated as in Section~2.2.
Final-evaluation network selection uses the already fixed threshold.

\begin{table}[H]
\caption{Common physical selection in execution order. Momenta and masses are in GeV; Appendix~\ref{app:fields} defines the branches.}
\label{tab:cuts}
\centering\small
\renewcommand{\arraystretch}{1.1}
\begin{tabularx}{\linewidth}{@{}>{\raggedright\arraybackslash}p{.23\linewidth}Y@{}}
\toprule
Stage & Implemented condition and purpose \\
\midrule
\textbf{Event trigger} & \texttt{trigE} or \texttt{trigM} is true, requiring an electron or muon trigger. \\
\midrule
\textbf{Lepton quality} & Apply the following requirements in order, retaining at least four leptons at each step, to select prompt-lepton candidates. \\
Flavour, identification & $|\texttt{lep\_type}|\in\{11,13\}$; \texttt{lep\_isTightID=true}. \\
Track, calorimeter isolation & $p_T>0$; sequentially require \texttt{lep\_ptcone30}/\texttt{lep\_pt}$<0.3$ and
\texttt{lep\_etcone20}/\texttt{lep\_pt}$<0.3$ to suppress objects with greater surrounding activity. \\
Transverse, longitudinal impact parameters & $|d_0/\sigma(d_0)|<$ 5 (electrons) or 3 (muons);
$|z_0/\cosh\eta|<0.5$ mm, limiting displacement relative to the primary vertex. \\
\midrule
\textbf{Four-lepton candidate} & The following sequential conditions define the common analysis population. \\
Multiplicity, trigger matching & Exactly four leptons after quality selection; at least one has \texttt{lep\_trigMatched=true}.
The original \texttt{lep\_n} may exceed four. \\
Ordered momenta, acceptance & In descending $p_T$ order, require $\geq20,15,10,7$;
electrons have $|\eta|<2.47$ and muons $|\eta|<2.7$, defining the kinematic and detector acceptance. \\
Charge, pairing & Zero total charge and two non-overlapping same-flavour opposite-sign (SFOS) pairs, forming a $ZZ^*$ candidate. \\
Pair masses & Every SFOS mass $>5$; $50<m_{Z1}<106$, $12<m_{Z2}<115$, excluding low-mass pairs and restricting the candidates. \\
Four-lepton mass, final state & $105\leq m_{4\ell}<140$; subsequently require two electrons and two muons, restricting the study population. \\
\bottomrule
\end{tabularx}
\end{table}

\subsection{Reference support and project differences}
The selection draws on published analyses without reproducing their complete procedures.
Sections 5.1--5.2 of the ATLAS 13 TeV analysis~\cite{ATLAS13} motivate the acceptance cuts, the leading three momentum thresholds, SFOS pairing, $Z$ windows and transverse impact-parameter cuts.
The subleading mass lower bound of $12\gev$ applies here because $m_{4\ell}<140\gev$.

The principal differences are:
\begin{itemize}[itemsep=1pt,topsep=2pt]
\item The reference uses loose electron identification and detailed muon requirements. This project applies tight-ID to both; working-point equivalence is unconfirmed.
\item Project momentum thresholds include equality and impose $p_T\geq7$ on the fourth lepton. The reference requires electron $E_T>7$ and muon $p_T>5$, with additional requirements for some muons.
\item Both project isolation cuts are $<30\%$. The reference uses track isolation $<15\%$, calorimeter isolation $<20\%$ for electrons and $<30\%$ for muons, with further cone and correction prescriptions.
\item Only exactly four qualified leptons in $2e2\mu$ are retained. The full multiple-candidate and angular-separation procedures are not reproduced. In this channel, the $Z$ windows are already stronger than the $5\gev$ pair cut.
\end{itemize}
The official educational example~\cite{HZZCuts,HZZAnalysis} supports some multiplicity, charge, flavour and momentum rules, but its isolation calls are commented out.
The Run-1 reference~\cite{ATLASRun1} supports the trigger-matching principle.
Exact external support for the tight-ID working point and $|z_0\sin\theta|<0.5$ mm remains unverified.
They remain fixed project choices and do not support a claim to reproduce a published measurement.
\FloatBarrier

\section{Error calculation and plot flags}\label{app:errors}

\subsection{Retention errors with shared events}
Selected events belong to the before-selection sample, so retention errors include covariance between the two yields.
For either process, let group $g$ contribute $T_g$ before selection and $H_g$ to the high-score category.
Then $T=\sum_gT_g$, $H=\sum_gH_g$ and $R=H/T$.
Under the independent-group Poisson-type convention used in the main text,
\[
 V_T=\sum_g T_g^2,\qquad V_H=\sum_g H_g^2,\qquad C_{HT}=\sum_g H_gT_g.
\]
$V_T$, $V_H$ are variance estimates and $C_{HT}$ the covariance estimate from shared group contributions.
For $T>0$, first-order error propagation gives
\[
 \sigma_R^2\simeq\frac{V_H}{T^2}+\frac{H^2V_T}{T^4}-\frac{2H C_{HT}}{T^3}.
\]
Taking the square root and multiplying by 100 gives the displayed retention error in percent.
Normalized spectra similarly include covariance between each bin and its total.
These formulas describe the convention used in the existing aggregates. They do not supply new error bars or significance tests for deletion differences.

\subsection{Statistical-support flags for mass bins}
Shading identifies bins whose MC estimates require caution.
Each bin is checked before selection and in the high-score category. It is shaded if either selection meets any of the following conditions:
\begin{enumerate}[leftmargin=1.7em,itemsep=2pt]
\item No physical event groups.
\item Nonpositive signed yield $Y$ or grouped variance $V$.
\item Signed effective count $N_{\rm eff}=Y^2/V<20$.
\item Cancellation ratio $\rho=|Y|/\sum_i|u_i|<0.2$.
\end{enumerate}
$u_i$ is the evaluation weight and $V$ the sum of squared group contributions in the bin.
$N_{\rm eff}$ checks weighted-sample support, while $\rho$ checks positive/negative cancellation.
A small $\rho$ means the signed yield is small relative to the absolute-weight sum.
Normalized shapes retain the same flags.
Shading denotes triggered checks, not an uncertainty interval or confidence band.

\subsection{Total-support conditions for normalized shapes}\label{app:shape-support}
Before returning normalized shapes, \texttt{mass\_spectra} checks the totals before and after selection separately.
For each curve, let $Y_{\rm total}$ be the total yield and $V_{\rm total}$ the variance after projection over event groups. These quantities and the absolute-weight sum define $N_{\rm eff}=Y_{\rm total}^2/V_{\rm total}$ and $\rho=|Y_{\rm total}|/\sum_i|u_i|$.
Both curves require positive total yield, $N_{\rm eff}\geq20$ and $\rho\geq0.2$.
Otherwise normalization is recorded as unsupported and the corresponding fractions are not returned.
The total variance includes covariance between bins from shared event groups. These conditions on the totals are separate from the checks used to shade individual bins.

\section{Implementation and reproduction}\label{app:network}

\subsection{Model and aggregate calculations}
\texttt{ResearchClassifier} is defined in \nolinkurl{src/higgsml/modeling/discriminants.py}; training is in \nolinkurl{src/higgsml/draft_v2/training.py}.
Hyperparameters are bound by the frozen protocol.
Training uses CPU float32. The loss function, \texttt{binary\_cross\_entropy\_with\_logits}, acts directly on the real-valued outputs to avoid numerical problems from applying sigmoid followed by logarithms.

Tables and figures use existing aggregates. Python's \texttt{json} reads the records; NumPy calculates ratios, within-seed contrasts and summaries; Matplotlib produces the figures.
Reproduction begins by checking the report, freeze and protocol identities and the completeness of the fifteen models. It then reads high- and low-score yields from \texttt{results.json} and checks their sum against the common yields before selection.
Retention and S/B use unrounded yields; within-seed contrasts precede medians and ranges.
Only presentation is rounded: four decimals for yields, two for percentages and percentage-point differences, and three for S/B and its differences.
Missing seeds and invalid metrics are flagged. Failed results are not removed before aggregation.
Build instructions are in \nolinkurl{paper/draft-v2/README.md}.

\subsection{Hidden-layer operations}\label{app:layers}
Fully connected layers combine previous-layer outputs; the first receives standardized inputs and subsequent layers receive hidden-layer outputs.
Training adjusts connection weights, biases and learnable normalization parameters.
LayerNorm normalizes values within the same event's hidden layer, then applies learnable scaling and shifting.
Its statistics differ from input standardization fitted across the training sample.
SiLU maps $a$ to $a/(1+e^{-a})$, providing nonlinearity; a stack of purely linear layers would reduce to one linear transformation.
During training, Dropout in the first two hidden layers sets outputs to zero with probability 0.1 and rescales the others by $1/0.9$. This reduces dependence on a few intermediate units.
Dropout is disabled for validation and final evaluation.

\subsection{Physical normalization settings}\label{app:normalization}
The configured $k$ factor and filter-efficiency multipliers are both 1; no additional factors are applied.
Production documentation is still needed to determine how branching fractions and filtering enter the effective cross section. This is necessary to avoid applying a factor twice.
The project does not apply the pileup, lepton and trigger scale factors used in the official example~\cite{HZZAnalysis}. Its yields therefore cannot be reproduced from the project settings alone.
These choices limit the physical interpretation of absolute yields.

\subsection{Threshold formula and illustrative example}\label{app:threshold}
Order the validation-background scores from lowest to highest. Let $w_i$ denote physical weights, $a_i=|w_i|$ their absolute values and $(i)$ the ordered indices.
Section~\ref{sec:threshold} implements
\[
 k=\min\left\{j:\sum_{i=1}^{j}a_{(i)}\geq\frac12\sum_i a_i\right\},
 \qquad t=s_{(k)}.
\]
The threshold is an observed score $s_{(k)}$, without interpolation, and ties are assigned together.
Cumulative weight can cross half the total, so even the validation high-score absolute-weight fraction need not be one half.
Validation and evaluation also contain different events, and yields use signed weights while the threshold uses absolute weights.
Together these explain why evaluation background retention need not equal 50\%.

\Needspace{9\baselineskip}
\textbf{Illustrative example (not actual run data).}
Suppose five validation-background scores are $0.10,0.30,0.60,0.80,0.90$, with absolute weights $1,2,1,4,2$.
The total is 10 and half is 5. Cumulative weights are $1,3,4,8,10$; the first to reach 5 has score 0.80, so $t=0.80$.
The $s\geq t$ selection accepts scores 0.80 and 0.90, with absolute-weight share $6/10=60\%$.
The example only illustrates the discrete rule and is not an actual run result.
\FloatBarrier

\Needspace{10\baselineskip}
\section{Result provenance and historical context}\label{app:provenance}

\subsection{Source of the retained results}
The numerical values are retained from the evidence record of 8 November 2026. It identifies the report path \nolinkurl{runs/h4l-draft-v2-draft001/report/}, report identity
\nolinkurl{a7e610a163e5d155705014201b7ee330323e3bf69309e0e89b2496edba5d4317}, and protocol \nolinkurl{config/protocols/h4l_draft_v2.json}.
The checks covered model and aggregate identities, file integrity, numerical consistency and document builds. Raw event payloads were not reaudited; see \nolinkurl{paper/draft-v2/BASIC-20261108-EVIDENCE.md}.
Literature sources and sample parameters follow the checks recorded on 30 October and 1 November 2026.
The numerical source identity remains unchanged. Revising the manuscript does not alter the frozen analysis or its exploratory qualification.
The CERN release restricts the collection to education and states that it is unsuitable for scientific publications. The report retains this educational and technical scope.

\subsection{What can be checked in the current working directory}\label{app:availability}
On 15 November 2026, the local directory with the recorded run name contained only preparation, training and freeze artifacts. The evaluation and report directories were absent.
The current prepared and freeze identities begin with \texttt{460ec2b5} and \texttt{a1303c3f}. They differ from the historical bindings and cannot replace the source of the reported numbers.

The audit checked controlled-MC file hashes and ROOT metadata. It also compared the current selection, inputs, allocation, weights, training, thresholds and spectrum-error formulas with the report, and ran the relevant software tests.
The fifteen metric rows and twelve deletion rows agree arithmetically with the existing evidence table.
This arithmetic check does not reconstruct all values from the missing unrounded results, events and model bindings.
The tables and figures retain their historical meaning. These checks do not constitute end-to-end numerical reproduction or independent physics validation.
Complete verification requires the exact result package and its upstream bindings. A matching directory name or numerically similar results are insufficient.

\subsection{Historical statistical methods}\label{app:interval}
Historical draft001 and validation001 studies examined interval widths using model-expected observations, profile-likelihood intervals, paired resampling and model-self pseudoexperiments.
Nominal undercoverage was observed, and independent coverage calibration was not established.
The retention and S/B calculations here do not use those historical results, and those diagnostics do not validate the current run.
Historical records must be distinguished by artifact identity because the same run name can refer to different artifacts.
Original numerical tables, identities and limitations remain in \nolinkurl{paper/draft-v2/DRAFT001-EVIDENCE.md} and its cited historical material.

\clearpage
\begingroup
\small
\renewcommand{\refname}{References}
\begin{thebibliography}{9}
\bibitem{ATLASSimulation}
ATLAS Collaboration, \emph{The ATLAS Simulation Infrastructure},
Eur.~Phys.~J.~C \textbf{70}, 823--874 (2010).
\href{https://doi.org/10.1140/epjc/s10052-010-1429-9}{doi:10.1140/epjc/s10052-010-1429-9};
\href{https://arxiv.org/abs/1005.4568}{arXiv:1005.4568}.
Reference for the general simulation chain, not proof of the specific production settings of these educational samples.

\bibitem{OpenData}
ATLAS Collaboration, \emph{ATLAS 13 TeV samples collection at least four leptons
(electron or muon), for 2020 Open Data release} (2020).
\href{https://opendata.cern.ch/record/15005}{CERN Open Data record 15005},
\href{https://doi.org/10.7483/OPENDATA.ATLAS.2Y1T.TLGL}{doi:10.7483/OPENDATA.ATLAS.2Y1T.TLGL}.
Source of the two files and loose release-preselection description; usage restrictions are in the record.

\bibitem{ATLAS13}
ATLAS Collaboration, \emph{Measurement of the Higgs boson coupling properties in
the $H\to ZZ^*\to4\ell$ decay channel at $\sqrt{s}=13$ TeV with the ATLAS detector},
JHEP \textbf{03} (2018) 095.
\href{https://doi.org/10.1007/JHEP03(2018)095}{doi:10.1007/JHEP03(2018)095};
\href{https://arxiv.org/pdf/1712.02304v2#page=8}{arXiv:1712.02304v2, Sections 5.1--5.2, pp. 8--9}.
Source for the cut and working-point comparison in Appendix~\ref{app:cuts}.

\bibitem{HZZCuts}
ATLAS Open Data outreach tools, \emph{HZZCuts.py}, revision
\nolinkurl{8ad2015be6d350060c3a183aff5802570bc9ada4}.
\href{https://github.com/atlas-outreach-data-tools/atlas-outreach-Python-uproot-framework-13tev/blob/8ad2015be6d350060c3a183aff5802570bc9ada4/HZZCuts.py#L2-L33}{Pinned source, lines 2--33}.
Multiplicity, charge, some momentum requirements, flavour and optional isolation definitions.

\bibitem{HZZAnalysis}
ATLAS Open Data outreach tools, \emph{HZZAnalysis.py}, same revision as Ref.~\cite{HZZCuts}.
\href{https://github.com/atlas-outreach-data-tools/atlas-outreach-Python-uproot-framework-13tev/blob/8ad2015be6d350060c3a183aff5802570bc9ada4/HZZAnalysis.py#L140-L168}{Pinned source, lines 140--168}.
Event-weight factors, applied cuts and commented-out isolation calls.

\bibitem{ATLASRun1}
ATLAS Collaboration, \emph{Measurements of Higgs boson production and couplings
in the four-lepton channel in $pp$ collisions at center-of-mass energies of
7 and 8 TeV with the ATLAS detector}, Phys.~Rev.~D \textbf{91}, 012006 (2015).
\href{https://doi.org/10.1103/PhysRevD.91.012006}{doi:10.1103/PhysRevD.91.012006};
\href{https://arxiv.org/pdf/1408.5191v3#page=7}{arXiv:1408.5191v3, Section 5.1, printed pp. 6--7}.
Reference for trigger matching; detailed thresholds and working points are not claimed equivalent to this project.

\bibitem{SampleInfo}
ATLAS Open Data outreach tools, \emph{infofile.py}, same revision as Ref.~\cite{HZZCuts}.
\href{https://github.com/atlas-outreach-data-tools/atlas-outreach-Python-uproot-framework-13tev/blob/8ad2015be6d350060c3a183aff5802570bc9ada4/infofile.py}{Pinned sample-parameter table}.
DSID, \texttt{xsec} and \texttt{sumw} for \texttt{ggH125\_ZZ4lep} and \texttt{llll}; not a complete generator/detector configuration record.
\end{thebibliography}
\endgroup
\end{document}
